\chapter{容器}
\label{ch:containers}
\chapterlead{容器是受隔离的进程，共享宿主机内核，因此启动快、开销小。本章用 podman 跑起一个服务，说明容器与 Nix 在职责上的分工，以及哪些场景不该用容器。}

\section{起步：运行第一个容器}

\begin{nixcode}
virtualisation.podman = {
  enable = true;
  dockerCompat = true; # 提供 docker 命令别名
};
\end{nixcode}

\begin{shellsession}
$ podman run --rm hello-world
$ podman run --rm -it alpine sh
$ podman ps -a
$ podman images
\end{shellsession}

\texttt{-{}-rm} 让容器退出后自动删除，\texttt{-it} 分配终端并保持交互。podman 默认以普通用户运行容器（rootless），比 docker 的默认行为安全。

\section{发布端口与挂载数据}

\begin{shellsession}
$ podman run -d --name web -p 8080:80 nginx:stable
$ curl -s localhost:8080 | head -3
$ podman logs -f web
$ podman stop web && podman rm web
\end{shellsession}

\texttt{-p 8080:80} 把宿主机的 8080 映射到容器的 80。\texttt{-d} 后台运行，\texttt{-{}-name} 起名字便于后续操作。数据默认随容器消失，需要保留时用卷：

\begin{shellsession}
$ podman volume create pgdata
$ podman run -d --name pg \
    -e POSTGRES_PASSWORD=secret \
    -v pgdata:/var/lib/postgresql/data \
    -p 5432:5432 postgres:16
\end{shellsession}

\section{容器与 Nix 的分工}

两者都解决“环境一致性”，但对象不同。

\begin{booktable}{\textwidth}{J{0.846} J{1.154}}
\toprule
维度 & 分工 \\
\midrule
\endhead
Nix & 构建期：可执行文件、库、配置文件从哪来 \\
容器 & 运行期：进程的隔离、资源限制、网络 \\
systemd & 生命周期：谁在什么时候启动、失败后怎么办 \\
\bottomrule
\end{booktable}

NixOS 上部署服务的自然顺序是：能用 systemd 服务直接跑的程序，就不放进容器；需要与宿主机隔离、或依赖一套完整运行时（例如某个特定版本的 Python 加一堆系统库）时，用容器；容器本身也由配置声明，而不是手工 \texttt{run}。

\begin{nixcode}
virtualisation.oci-containers = {
  backend = "podman";
  containers.caddy = {
    image = "docker.io/library/caddy:2";
    ports = [
      "80:80"
      "443:443"
    ];
    volumes = [ "/var/lib/caddy:/data" ];
  };
};
\end{nixcode}

这段配置让容器随系统重建启动，状态由 systemd 管理。手工 \texttt{podman run} 的容器在重建后不会自动恢复，因此只适合实验。

\section{自己构建镜像}

\begin{textcode}
# Containerfile
FROM alpine:3.20
RUN apk add --no-cache curl
COPY entry.sh /entry.sh
ENTRYPOINT ["/entry.sh"]
\end{textcode}

\begin{shellsession}
$ podman build -t myscript:0.1 .
$ podman run --rm myscript:0.1
\end{shellsession}

在 NixOS 上还有另一条路：用 \texttt{dockerTools.\allowbreak{}buildImage} 在 Nix 里构造镜像，内容和依赖都来自 store，可复现性与系统一致。

\begin{nixcode}
pkgs.dockerTools.buildImage {
  name = "myscript";
  tag = "0.1";
  copyToRoot = pkgs.buildEnv {
    name = "image-root";
    paths = [
      pkgs.bash
      pkgs.curl
    ];
  };
  config.Cmd = [ "/bin/bash" ];
}
\end{nixcode}

\section{什么时候不该用容器}

三类情况不适合容器。第一，只有一个静态二进制的小服务，直接用 systemd 更简单。第二，需要直接访问硬件（GPU、USB 设备）的工具，容器只增加一层配置。第三，把容器当作“隔离安全边界”：容器共享内核，隔离强度弱于虚拟机，运行不可信代码应使用虚拟机。

\begin{notebox}[rootless 的限制]
普通用户运行容器时，无法绑定 1024 以下的端口，也不能挂载任意路径。这两个限制都可以绕开，但需要额外配置。服务器上部署对外服务时，通常在 80 与 443 前面放一个反向代理，容器监听高位端口即可。
\end{notebox}

\section{练习}

\begin{enumerate}
  \item 运行一个 nginx 容器并映射到 8080，用 \texttt{curl} 验证，再查看日志。
  \item 用卷运行 PostgreSQL，重启容器后确认数据仍在。
  \item 写一个 Containerfile，构建并运行一个打印时间的镜像。
  \item 用 \texttt{dockerTools.\allowbreak{}buildImage} 构造同样功能的镜像，比较两者大小。
  \item 判断本机上三个现有服务，哪些适合放进容器，说明理由。
\end{enumerate}

\begin{checkpoint}
本章结束时，应能：
\begin{itemize}
  \item 用 podman 运行、发布端口、查看日志、清理容器；
  \item 用卷保存容器数据，说明容器删除后什么会丢；
  \item 说清 Nix、容器、systemd 三者的职责分工；
  \item 用 \texttt{virtualisation.oci-containers} 声明式地运行容器；
  \item 判断哪些服务不该容器化。
\end{itemize}
\end{checkpoint}
